\documentclass[12pt,a4paper]{article}
\usepackage[utf8]{inputenc}
\usepackage[T1]{fontenc}
\usepackage[french]{babel}
\usepackage{amsmath}
\usepackage{geometry}
\usepackage{setspace}
\usepackage{hyperref}
\geometry{margin=2.5cm}
\setstretch{1.25}

\title{Revue de Synthèse : Mobilité, salaires et trajectoires professionnelles en France}
\author{Thomas Barat, Chloé Martin \& Ilayda Yilmaz}
\date{Octobre 2025}

\begin{document}
\maketitle

\section*{Introduction}

def utilité 
def théorie capital humain
def théorie de l'appariement
def modèle de Mincer
def mobilité (effets)
    - inter (Simonnet, Kramatz)
    - intra
    - sectorielle
    - géographique (Magrier)

Kramaz : 3 phases 
- MCO classique
- économétrie des panels
- modèle structurel (def)

\section*{Données}

Présentation des variables récurrentes


\section*{Présentation de notre modèle}

Modèle Maigrier : modèle à équations simultanées permettant d’étudier à la fois l’impact de diverses variables, dont le salaire, sur la probabilité de migrer et l’impact de la migration sur les salaires.
Modèle Simmonet : trouver des articles pour la justification de l'ajout de la variable ancienneté
modèle d'Heckman : traitement du biais de sélection dans les 
décomposition
lawel

\section*{Limites}

Limites de nos articles

\end{document}